\chapter{Natural Harmony as Informational Efficiency}

\epigraph{\emph{Recurrence conserves intelligibility.}}{}

The previous chapters established three conditions for the present argument. Harmonic information was defined in relational rather than substantial terms, consonance was reformulated as a situated achievement of coordinated stability, and the program's architecture of evidence made clear that strong claims require explicit warrants rather than intuitive continuity alone. Chapter 8 begins from that prepared ground. Its task is to ask whether natural harmonic organization can be understood not only as recurrent structure, but as a candidate regime of informational efficiency.

That question must be stated with care at the outset, because the direction of the argument remains open. The chapter does not assume that harmony has already been proven to cause efficiency, nor that efficiency has already been shown to generate harmonic organization. The working hypothesis is narrower: under conditions where systems must coordinate, stabilize, interpret, or predict patterned relations, some harmonic regimes may reduce corrective burden relative to nearby alternatives. Whether harmonic organization drives that economy, whether pressures for economy recruit harmonic organization, or whether both co-emerge under shared constraints remains an open problem. The chapter therefore advances a directional hypothesis, not a closed mechanism.

\section{From ratio to recurrence}

The first step in the argument is formal and comparatively modest. If two oscillatory components stand in the relation $r = f_2 / f_1$, then the question is whether some values of $r$ are more likely than others to generate recurrent organization. HIT begins from the observation that simple proportional relations often yield more legible periodic return, more stable interference structure, and wider basins of coordinated behavior than more complex or incommensurable alternatives. This is visible in multiple empirical and formal archives at once: in Lissajous figures, in recurrence analysis, in mode-locking, and in coupled-oscillator dynamics. None of these literatures alone establishes the larger theory, but together they justify treating recurrence as the first operational bridge between harmonic simplicity and efficiency.

Kuramoto-style models make the point in a language of coupled phases rather than musical intervals. In their standard form, they describe the dynamics of interacting oscillators as:

\begin{equation}
\frac{d\theta_i}{dt} = \omega_i + \frac{K}{N} \sum_j \sin(\theta_j - \theta_i)
\end{equation}

where $\omega_i$ denotes the natural frequency of oscillator $i$, $K$ the coupling strength, and the sine term the phase interaction among oscillators \parencite{kuramoto_1984}. What matters for HIT is not the full mathematical machinery of the model, but the general lesson that some relational regimes support more robust locking than others. Rational relations with smaller integer structure often permit more stable coordination, wider tongues of capture, and more persistent patterned behavior than highly complex or effectively incommensurable relations do \parencite{acebron_2005, cartwright_2001}. This does not show that simple ratios are efficient in themselves. It shows that they are recurrently available to systems seeking coordinated persistence.

Recurrence quantification analysis provides a second formulation of the same tendency in state-space terms. A recurrence matrix can be written as:

\begin{equation}
R_{ij} = \Theta(\varepsilon - \| \mathbf{x}_i - \mathbf{x}_j \|)
\end{equation}

where $\Theta$ is the Heaviside step function and $\varepsilon$ sets the neighborhood threshold. A minimal recurrence-rate measure then takes the form:

\begin{equation}
RR = \frac{1}{N^2} \sum_i \sum_j R_{ij}
\end{equation}

The important point is conceptual before it is technical. Recurrence measures do not ask whether a system looks numerically elegant in the abstract; they ask how often trajectories return near prior states and how legible that return becomes over time. Trulla and colleagues showed that simpler intervallic relations in auditory contexts correlate with stronger and cleaner recurrence structure than more complex alternatives \parencite{trulla_2018}. Gallozzi and Strollo show a related tendency in Lissajous configurations, where simple relations tend toward closure and topological economy, whereas more complex relations produce increasingly elaborate and less recurrent figures \parencite{gallozzi_2023}. In driven media, pattern formation under periodic forcing likewise demonstrates that some relational regimes stabilize faster and more legibly than others, even when the medium itself remains noisy or nonlinear \parencite{engels_2007, nguyen_2019}.

At this point, however, the argument must introduce a distinction it will need for everything that follows. Not every recurrence is informative. A trivially repetitive regime may be perfectly regular while remaining functionally empty for the system under consideration. What matters to HIT is not recurrence as mere return, but recurrence that preserves discriminability, supports uptake, and stabilizes a difference capable of making a difference. The question is not whether $2{:}1$ returns more neatly than $\sqrt{2}{:}1$ in a mathematical sense. It is whether some recurrent relations organize the field in ways that another system can use, track, or remain coupled to. That narrower notion of informative recurrence is the real bridge to efficiency.

\section{Informative recurrence and reduced corrective burden}

Recurrence becomes theoretically important for HIT when it is read not only as pattern, but as a constraint on the amount of correction a system must perform in order to remain organized. A highly recurrent regime does not need to invent its order afresh at every moment. It can return to nearby states, exploit already available regularities, and stabilize expectation over time. Yet the relevant regime is not one of perfect rigidity. A system with no variation left to register would be orderly in a dead sense: maximally repetitive and minimally informative. The chapter's target is different. It is the intermediate regime in which there is enough recurrence to economize correction and enough variation to preserve discriminability, responsiveness, and use. In that sense the phrase \enquote{organized chaos} names a real theoretical problem rather than a literary flourish.

Landauer's principle provides a minimal lower-bound intuition for why this matters. In its most famous form, it states that erasing one bit of information requires at least $kT \ln 2$ of dissipated energy \parencite{landauer_1961}. Bennett's later work clarified the same terrain by showing that reversible computation can in principle eliminate thermodynamic dissipation entirely, localizing the unavoidable cost specifically to logically irreversible operations such as erasure \parencite{bennett_1982}. These results do not demonstrate that harmonic relations are cheap in organisms, brains, or ecosystems. They do establish that information processing is not free in principle and that any theory of organization must eventually face the cost of state transformation. The word entropy must therefore be handled carefully in what follows. Sometimes it refers to thermodynamic cost in the strict sense that makes Landauer's bound meaningful; sometimes HIT uses it in a broader, explicitly reformulated sense to name changes in order, legibility, and constraint across domains. These senses are related but not interchangeable, and neither should be collapsed into Shannon's narrower question of uncertainty reduction, even when the three terrains partially overlap.

Still and colleagues provide the bridge closest to HIT's actual concern by connecting predictive structure to thermodynamic efficiency. Their work suggests that systems able to capture regularities in their environments can reduce wasted dissipation relative to systems burdened by weakly structured or weakly predictable inputs \parencite{still_2012}. Friston's free-energy framework points in a related direction: systems persist by minimizing the discrepancy between expected and incoming states, thereby lowering the corrective work needed to maintain viable organization \parencite{friston_2010}. The present claim remains conditional. Where systems must track, predict, integrate, or stabilize patterned relations, informative recurrence often implies less corrective work than weaker or noisier organization. Harmonic simplicity matters here because it is one common generator of such recurrence, not because numerical simplicity by itself settles the argument.

Recent work on probabilistic hardware architectures sharpens the same problem from a different engineering frontier. Jelin{\v{c}}i{\v{c}} and colleagues describe a system in which the cost of sampling depends not only on how expressive a target distribution is, but on the geometry of the energy landscape itself: basin depth, barrier height, and traversability all become computationally decisive \parencite{jelincic_2025}. In the Boltzmann formalism underneath that architecture, what matters are relative probabilities among states, $P(x_1)/P(x_2)=e^{-(E_1-E_2)/kT}$, rather than isolated states considered one by one. Read at that scale, the paper reinforces a point central to HIT: organization can alter cost because access is governed by relations within a field rather than by the isolated inventory of possible states.

The same paper suggests a second bridge best stated as a conjecture. Its diffusion-thermodynamic chain decomposes a multimodal landscape into a sequence of locally simpler denoising steps. The harmonic series likewise presents a complex organization through a hierarchy in which each level refines the previous one through comparatively simple local constraints. The two constructions belong to different formalisms, but the structural resemblance is worth naming. If that analogy later proves formalizable, it would strengthen the idea that some efficiencies arise by staging complexity rather than by flattening it.

This makes the chapter's chain more exact than it first appears. Simple ratios tend, under suitable conditions, to generate recurrent structure. Only some of that recurrence is informative in the relevant sense. Where it is informative, it can reduce the need for continual large-scale correction. Reduced corrective burden is one defensible sense of informational efficiency. That is a stronger claim than saying that periodic return looks elegant, and a weaker one than saying that harmony has already been demonstrated as the universal optimum of natural systems.

\section{Informational efficiency and living organization}

The argument becomes more consequential when it moves from abstract dynamics to living organization. Biological systems do not merely oscillate. They must coordinate under noisy conditions, preserve viability, and respond without disintegrating. In such settings, efficiency is never just a matter of saving energy in the narrow engineering sense. It is also a matter of preserving coherent organization while minimizing destabilizing correction. Three readings therefore need to remain distinct. One is that harmonic organization helps produce efficiency. Another is that efficient systems tend to recruit harmonic organization. A third is that both arise together under shared dynamic constraints. The present chapter can motivate all three possibilities, but it cannot yet decide among them.

Neural dynamics offer one of the clearest domains in which this ambiguity becomes productive rather than paralyzing. Canolty and Knight showed that cross-frequency coupling is not random clutter layered over neural activity, but a structured mechanism through which different temporal scales can be coordinated within the brain \parencite{canolty_2010}. Jacobs and colleagues, working on traveling waves and cortical geometry, suggest a related lesson: local and global coordination often depend on orderly relations among oscillatory processes rather than on isolated point events alone \parencite{jacobs_2025}. These results support the idea that ratio-structured coordination can accompany efficient integration across scales. They do not by themselves establish whether harmonicity is the cause, the consequence, or one manifestation of a broader economy of organization.

Physiological oscillators point in the same direction. Glass and Mackey's work on biological rhythms demonstrates that viable living systems often preserve order not through rigid uniformity, but through flexible locking among interacting cycles \parencite{glass_1988}. Cardiorespiratory coordination, endocrine timing, and other rhythmic regimes do not remain healthy by eliminating difference, but by stabilizing it into low-conflict relation. Here again the interesting regime is neither perfect regularity nor uncontrolled noise. It is an intermediate one: enough recurrent structure to reduce destabilizing correction, enough plasticity to remain adaptive. Harmonic relations become theoretically interesting in that context because they may mark families of regimes within which such economy is easier to sustain.

The ecological case is similar. Krause's work on the acoustic niche shows that biophonic organization is not a chaotic pile-up of signals, but often a partitioned spectral and temporal arrangement in which organisms reduce destructive overlap while preserving communicative salience \parencite{krause_2013}. An ecology that lowers masking while preserving legible difference has solved a problem of coordination under constraint. Whether one interprets that as harmony producing efficiency or efficiency selecting harmonic spacing is precisely the kind of directionality question the chapter must leave open. The literature is enough to justify the convergence. It is not yet enough to close the mechanism.

An autopoietic reading can sharpen the stakes of that convergence without pretending to settle it. If a living system is organizationally precarious, then efficiency matters because organization can fail. Di Paolo's extension of autopoiesis is helpful here: autonomy becomes normatively charged precisely because the system must continually preserve itself under conditions that can dissolve its viability \parencite{dipaolo_2005}. In that light, the chain developed in this chapter can be reread more cautiously. Ratios that support informative recurrence may reduce corrective burden not only in an abstract informational sense, but because they help sustain organization under conditions of vulnerability. This remains an interpretive bridge, not a formal derivation. Autopoiesis still lacks a metric directly comparable to recurrence quantification or variational free energy. But the bridge is conceptually important because it translates efficiency into the biological language of maintaining organization under precarious conditions.

\section{The biosphere and semiosphere of recurrence}

Chapter 8 becomes most distinctive when recurrence is read not only as dynamical stabilizer, but also as semiotic economy. A recurrent pattern matters because it can be recognized again. Recognition lowers interpretive burden. It allows a receiving system to anticipate, classify, remember, and coordinate without treating each event as wholly novel. In this sense, recurrence does not merely conserve motion. It conserves intelligibility. The same structured return that lowers corrective burden at the level of dynamics can lower ambiguity at the level of interpretation.

Bateson's formulation is decisive here and should enter before the argument drifts into thermodynamic abstraction. A relation becomes informative when it changes what another system can do, anticipate, or sustain. From that perspective, the distinction introduced earlier between trivial and informative recurrence acquires its full force. A perfectly repetitive regime with no differential consequence may be orderly but informationally thin. By contrast, a regime that returns with enough stability to be expected and enough variation to remain discriminable becomes semiotically economical. It can be taken up without being reduced to noise, and it can preserve difference without requiring total reinterpretation at every step.

This is why recurrence can be understood as a compression device across levels of organization. A recurrent structure is easier to store, easier to detect, and easier to coordinate around than a continually novel one, provided that the receiving system is tuned to the relevant regularity. Compression here should not be heard only in the narrow algorithmic sense. It also names a reduction in interpretive labor. Fewer ambiguities need to be resolved, fewer possible states need to be entertained at once, and fewer contextual repairs are needed to keep the signal usable. In that sense, the same regime can operate simultaneously as dynamical economy and semiotic economy.

Krause's acoustic ecologies can be reread fruitfully through this lens. Spectral partitioning does not merely reduce collision; it also makes signals more discriminable within a crowded field. A bird call, insect rhythm, or environmental signal becomes easier to separate, track, and respond to when the field itself is recurrently organized rather than maximally congested. The same logic applies more broadly wherever repeated structure coordinates action across scales: lower overlap, more stable cues, fewer interpretive dead ends. HIT does not need to claim that all biospheric order is harmonic in a strict ratio-theoretic sense. It needs only the more disciplined point that recurrent relational pattern can serve at once as a stabilizer of organization and as an economy of meaning.

The ecological and biosemiotic extension should therefore end on restraint rather than inflation. Natural harmony is not being elevated into a total explanation of life, nor into a secret code of the biosphere. It is being proposed as one family of recurrent organizations that may reduce both corrective and interpretive burden under conditions of complexity. That proposal is narrower than metaphysics and stronger than metaphor. It is also the point at which HIT differs most clearly from a merely acoustic or purely physical theory of order.

\section{A minimal proposition for HIT}

The chapter can now compress its argument into a minimal proposition. Under suitable material and interpretive conditions, simpler harmonic relations tend to generate more recurrent organization than more complex or weakly commensurable alternatives. Where that recurrence remains informative rather than trivial, it can reduce the need for continual large-scale correction in systems that must stabilize, predict, or coordinate patterned behavior. When that happens, natural harmony may be treated as an attractor of informational efficiency. The phrase does not mean that harmony is free, perfect, or universally privileged. It means that some harmonic regimes can support coherent organization at lower corrective and interpretive cost than nearby alternatives.

This proposition remains explicitly tiered. At the level of observation, the literature supports the association between harmonic simplicity and stronger recurrence in multiple domains. At the level of hypothesis, HIT proposes that informative recurrence can often imply reduced corrective burden. At the level of inference, the chapter argues that natural harmony may therefore be understood as a candidate regime of informational efficiency. The direction of causality remains open: efficiency may recruit harmonic organization, harmonic organization may support efficiency, or both may emerge together under shared constraints. The steps are linked but not collapsed.

The leanest formulation is therefore also the safest one. Natural harmony is theoretically important not only because it is recurrent, but because informative recurrence may be one of the most basic ways in which systems economize the work of staying organized while preserving legible difference. If that is right, then harmony is not just pleasant structure, not just stable ratio, and not just cultural artifact. It becomes one candidate answer to a deeper question: how does a system hold coherence without paying the highest possible cost for every correction?

The next step follows directly. If some forms of harmonic organization function as efficiency attractors, then a biological question immediately arises: can living systems detect, prefer, or orient toward such regimes in their own operation? Chapter 9 takes up that question by asking how consonance becomes not only a formal property, but a sense.
